(1) 微分方程一般解为 \(B_z=C\cosh(x/\lambda_L)+D\sinh(x/\lambda_L)\)。边界 \(B_z(a)=B_z(-a)=B_0\) 导致 \(D=0\)、\(C=B_0/\cosh(a/\lambda_L)\)，故 \(B_z(x)=B_0\cosh(x/\lambda_L)/\cosh(a/\lambda_L)\)，在中心最弱、两侧为 \(B_0\)。
(2) 仅 \(B_z(x)\) 非零时 \(\nabla\times\mathbf B=-B_z'(x)\hat y\)，因此 \(J_y(x)=-B_0\sinh(x/\lambda_L)/[\mu_0\lambda_L\cosh(a/\lambda_L)]\)，关于中心为奇函数，两侧屏蔽电流方向相反。若 \(a\gg\lambda_L\)，磁场只在离两表面数个 \(\lambda_L\) 的区域显著，表现为 Meissner 屏蔽。
(3) 单位面积总储能 \(U/A=\int_{-a}^a[B_z^2/(2\mu_0)+\mu_0\lambda_L^2J_y^2/2]\,dx\)。代入双曲函数，括号为 \(B_0^2[\cosh^2(x/\lambda_L)+\sinh^2(x/\lambda_L)]/[2\mu_0\cosh^2(a/\lambda_L)]\)；积分得 \(U/A=B_0^2\lambda_L\tanh(a/\lambda_L)/\mu_0\)。
(4) 若 \(a/\lambda_L\ll1\)，则 \(\tanh(a/\lambda_L)\sim a/\lambda_L\)，\(U/A\sim B_0^2a/\mu_0\)，等于近似均匀场填充厚度 \(2a\) 的纯磁能。若 \(a/\lambda_L\gg1\)，\(U/A\to B_0^2\lambda_L/\mu_0\)，能量只储存在两侧穿透层。\([B'/\mu_0]=\mathrm{A\,m^{-2}}\)，与体电流密度相符。
